\renewcommand{\thechapter}{\arabic{chapter}}

\chapter{ER diagram relační databáze}
\label{app:er-diagram}

\begin{sloppypar}
Diagram znázorňuje strukturu relační databáze PostgreSQL. Společné sloupce \texttt{created\_at} a~\texttt{updated\_at} jsou pro přehlednost vynechány.
\end{sloppypar}

\begin{figure}[H]
\centering
\resizebox{\textwidth}{!}{%
\begin{tikzpicture}[
    head/.style={
        rectangle, draw, fill=blue!12,
        minimum width=3.2cm, minimum height=0.55cm,
        font=\sffamily\bfseries\footnotesize
    },
    body/.style={
        rectangle, draw, fill=white,
        minimum width=3.2cm,
        font=\ttfamily\tiny,
        text width=2.9cm,
        inner sep=2.5pt,
        align=left,
        anchor=north
    },
    rel/.style={-{Stealth[length=2mm]}, semithick, black!50},
    rellbl/.style={font=\sffamily\tiny\bfseries, fill=white, inner sep=1pt}
]

% === SLOUPEC 1 ===

\node[head] (uh) at (0, 0) {users};
\node[body, below=-0.3pt of uh] (ub) {%
\textbf{PK} id: uuid\\
email: varchar\\
password\_hash: text\\
name: varchar\\
role: varchar
};

\node[head] (eh) at (0, -4.2) {events};
\node[body, below=-0.3pt of eh] (eb) {%
\textbf{PK} id: uuid\\
title, slug: varchar\\
description: text\\
location: varchar\\
start\_date, end\_date\\
cover\_image: text\\
event\_type: varchar\\
special\_status: varchar\\
published: boolean
};

\node[head] (rsh) at (0, -8.8) {recruitment\_submissions};
\node[body, below=-0.3pt of rsh] (rsb) {%
\textbf{PK} id: uuid\\
name, email, phone\\
education, experience\\
interests, motivation\\
skills, preferred\_role\\
availability, task: text\\
status: varchar\\
notes: text
};

% === SLOUPEC 2 ===

\node[head] (ah) at (4.5, 0) {articles};
\node[body, below=-0.3pt of ah] (ab) {%
\textbf{PK} id: uuid\\
title, slug: varchar\\
excerpt, content: text\\
cover\_image: text\\
images, tags: jsonb\\
category, status: varchar\\
published: boolean\\
published\_at: timestamp\\
\textbf{FK} author\_id → users
};

\node[head] (mh) at (4.5, -4.2) {members};
\node[body, below=-0.3pt of mh] (mb) {%
\textbf{PK} id: uuid\\
name, email, phone\\
role, department\\
tags: jsonb\\
bio, avatar: text\\
linkedin, github: text\\
active: boolean
};

\node[head] (seth) at (4.5, -8.8) {recruitment\_settings};
\node[body, below=-0.3pt of seth] (setb) {%
\textbf{PK} id: uuid\\
is\_active: boolean\\
roles: jsonb\\
tasks: jsonb
};

% === SLOUPEC 3 ===

\node[head] (ph) at (9, 0) {projects};
\node[body, below=-0.3pt of ph] (pb) {%
\textbf{PK} id: uuid\\
name, slug: varchar\\
description: text\\
category, status: varchar\\
cover\_image: text\\
images, specs: jsonb\\
published: boolean\\
is\_featured: boolean
};

\node[head] (pmh) at (9, -4.2) {project\_members};
\node[body, below=-0.3pt of pmh] (pmb) {%
\textbf{PK} id: uuid\\
\textbf{FK} project\_id → projects\\
\textbf{FK} member\_id → members\\
role: varchar
};

\node[head] (parh) at (9, -8.8) {partners};
\node[body, below=-0.3pt of parh] (parb) {%
\textbf{PK} id: uuid\\
name: varchar\\
tier: enum (s/g/d)\\
logo: text\\
description: text\\
full\_description: text\\
hero\_image: text\\
website: text\\
order: integer\\
published: boolean
};

% === RELACE ===

% users 1:N articles
\draw[rel] (uh.east) -- node[rellbl, above] {1:N} (ah.west);

% projects 1:N project_members
\draw[rel] (ph.south west) -- node[rellbl, left] {1:N} (pmh.north west);

% members 1:N project_members
\draw[rel] (mh.east) -- node[rellbl, above] {1:N} (pmh.west);

\end{tikzpicture}
}%
\caption{ER diagram relační databáze platformy.}
\label{fig:er-diagram}
\end{figure}

\chapter{Screenshoty aplikace}
\label{app:screenshots}

V~této příloze jsou souhrnně zobrazeny vybrané screenshoty aplikace ilustrující klíčové části platformy. Obrázky jsou rovněž průběžně odkazovány v~jednotlivých kapitolách u~odpovídajících popisů.

\begin{figure}[H]
\centering
\includegraphics[width=\textwidth]{homepage.png}
\caption{Úvodní stránka veřejného webu Czech Rocket Society.}
\label{fig:appendix-homepage}
\end{figure}

\begin{figure}[H]
\centering
\includegraphics[width=\textwidth]{recruitment-form.png}
\caption{Náborový formulář -- ukázka jednoho z~kroků s~vizualizací postupu.}
\label{fig:appendix-recruitment}
\end{figure}

\begin{figure}[H]
\centering
\includegraphics[width=\textwidth]{admin-events.png}
\caption{Administrační rozhraní -- správa událostí s~přehledem typů, statusů a~viditelnosti.}
\label{fig:appendix-admin-events}
\end{figure}

\begin{figure}[H]
\centering
\includegraphics[width=\textwidth]{richtext-editor.png}
\caption{Rich-text editor Tiptap v~administraci s~formátovacím toolbarem.}
\label{fig:appendix-richtext}
\end{figure}

\begin{figure}[H]
\centering
\includegraphics[width=\textwidth]{telemetry-live.png}
\caption{Telemetrický dashboard v~live režimu -- časové grafy, senzory, GPS trajektorie a~stavová lišta.}
\label{fig:appendix-telemetry}
\end{figure}

\chapter{Obsah přiloženého média}
\label{app:medium}

% TODO: Doplnit odkaz na GitHub repozitář

\noindent Repozitář obsahuje:
\begin{itemize}
    \item \texttt{apps/web/} -- zdrojový kód frontendové aplikace (Next.js)
    \item \texttt{apps/api/} -- zdrojový kód backendové aplikace (Hono)
    \item \texttt{apps/telemetry/} -- zdrojový kód telemetrického serveru (Go)
    \item \texttt{docker/} -- konfigurace Docker Compose a~služeb
    \item \texttt{thesis/} -- zdrojové soubory bakalářské práce (\LaTeX)
    \item \texttt{README.md} -- návod k~instalaci a~spuštění
\end{itemize}

\chapter{Zadání práce}
\label{app:zadani}

\noindent Originální zadání bakalářské práce z~informačního systému IS~STAG je vloženo na~následujících stránkách.

\includepdf[pages=-,pagecommand={},width=\textwidth]{images/zadani-stag.pdf}
